% GENERATED FILE. Edit canonical lessons, then run npm run generate:books.
% canonical-source-hash: fnv1a64:606944f4d1d7c328
% canonical-lessons: IT-C445-campanile, IT-C445-duomo, IT-C445-cattedrale, IT-C445-portico, IT-C445-loggia

\chapter{In piazza}
\label{ch:it-in-piazza}
\hlchaptermodality{445}

\begin{chapteropening}
\textbf{By the end of this chapter:} \emph{I can describe the buildings round a square: the bell tower, the main church, the cathedral, the arcades and the loggia.}

Complete the last lesson of chapter 445: I can describe the buildings round a square: the bell tower, the main church, the cathedral, the arcades and the loggia.
\end{chapteropening}

\section[il campanile]{il campanile — the bell tower}

\label{lesson:IT-C445-campanile}



\begin{quote}
Before the new one: say the Italian for the house number, then the Italian for the intercom.
\end{quote}

\subsection*{You'll want to know: il campanile}
\textbf{il campanile} — \textquotedblleft{}the bell tower\textquotedblright{}.

\begin{quote}
Dal campanile si vede tutta la città.
\end{quote}

\begin{quote}
\emph{From the bell tower the whole city is in view.}
\end{quote}

\emph{Il campanilismo} — loyalty to your own town and its bell tower.

\begin{cousinweb}
\emph{campanile} is built on \emph{campana} 'bell', which is inherited from Late Latin \emph{campāna} (EtymDB).
\end{cousinweb}

\begin{grammarlens}[title={What you've built}]
One more word for telling, explaining and describing at B1.
\end{grammarlens}

\subsection*{Your turn}
\begin{itemize}
\raggedright
  \item \emph{Say it:} \emph{il campanile}
  \item \emph{Say it:} \emph{il campanile}, once more
  \item \emph{Say it:} a sentence of your own with \emph{il campanile}
  \item \emph{Recall:} say the Italian for the house number, then the Italian for the intercom, then say \emph{il campanile} again
\end{itemize}

\begin{tcolorbox}[breakable,colback=teal!4,colframe=teal!35!black,title={Before you move on}]
What does il campanile mean? (\textquotedblleft{}The bell tower\textquotedblright{}.) Say it once more.
\end{tcolorbox}
\section[il duomo]{il duomo — the main church (of a town), the cathedral}

\label{lesson:IT-C445-duomo}



\begin{quote}
Before the new one: say the Italian for the intercom, then the Italian for the bell tower.
\end{quote}

\subsection*{You'll want to know: il duomo}
\textbf{il duomo} — \textquotedblleft{}the main church (of a town), the cathedral\textquotedblright{}.

\begin{quote}
Il duomo è in fondo alla via.
\end{quote}

\begin{quote}
\emph{The main church is at the end of the street.}
\end{quote}

Many towns call their main church \emph{il Duomo}.

\begin{cousinweb}
\emph{duomo} is inherited from Latin \emph{domus} 'house', short for \emph{domus ecclēsiae} 'the house of the church' (EtymDB).
\end{cousinweb}

\begin{grammarlens}[title={What you've built}]
One more word for telling, explaining and describing at B1.
\end{grammarlens}

\subsection*{Your turn}
\begin{itemize}
\raggedright
  \item \emph{Say it:} \emph{il duomo}
  \item \emph{Say it:} \emph{il duomo}, once more
  \item \emph{Say it:} a sentence of your own with \emph{il duomo}
  \item \emph{Recall:} say the Italian for the intercom, then the Italian for the bell tower, then say \emph{il duomo} again
\end{itemize}

\begin{tcolorbox}[breakable,colback=teal!4,colframe=teal!35!black,title={Before you move on}]
What does il duomo mean? (\textquotedblleft{}The main church (of a town), the cathedral\textquotedblright{}.) Say it once more.
\end{tcolorbox}
\section[la cattedrale]{la cattedrale — the cathedral}

\label{lesson:IT-C445-cattedrale}



\begin{quote}
Before the new one: say the Italian for the bell tower, then the Italian for the main church.
\end{quote}

\subsection*{You'll want to know: la cattedrale}
\textbf{la cattedrale} — \textquotedblleft{}the cathedral\textquotedblright{}.

\begin{quote}
La cattedrale ha tre porte di bronzo.
\end{quote}

\begin{quote}
\emph{The cathedral has three bronze doors.}
\end{quote}

The bishop's church; \emph{la cattedra} is the bishop's chair, and also a teacher's desk.

\begin{cousinweb}
\emph{cattedrale} is a learned borrowing from Late Latin \emph{cathedrālis}, from \emph{cathedra} 'chair', from Greek \emph{kathédra} 'seat' (EtymDB): the church with the bishop's chair.
\end{cousinweb}

\begin{grammarlens}[title={What you've built}]
One more word for telling, explaining and describing at B1.
\end{grammarlens}

\subsection*{Your turn}
\begin{itemize}
\raggedright
  \item \emph{Say it:} \emph{la cattedrale}
  \item \emph{Say it:} \emph{la cattedrale}, once more
  \item \emph{Say it:} a sentence of your own with \emph{la cattedrale}
  \item \emph{Recall:} say the Italian for the bell tower, then the Italian for the main church, then say \emph{la cattedrale} again
\end{itemize}

\begin{tcolorbox}[breakable,colback=teal!4,colframe=teal!35!black,title={Before you move on}]
What does la cattedrale mean? (\textquotedblleft{}The cathedral\textquotedblright{}.) Say it once more.
\end{tcolorbox}
\section[il portico]{il portico — the arcade, the portico}

\label{lesson:IT-C445-portico}



\begin{quote}
Before the new one: say the Italian for the main church, then the Italian for the cathedral.
\end{quote}

\subsection*{You'll want to know: il portico}
\textbf{il portico} — \textquotedblleft{}the arcade, the portico\textquotedblright{}.

\begin{quote}
Quando piove camminiamo sotto i portici.
\end{quote}

\begin{quote}
\emph{When it rains we walk under the arcades.}
\end{quote}

Often plural: \emph{i portici}.

\begin{cousinweb}
\emph{portico} is a learned borrowing from Latin \emph{porticus} 'colonnade, porch' (EtymDB), from \emph{porta} 'gate'.
\end{cousinweb}

\begin{grammarlens}[title={What you've built}]
One more word for telling, explaining and describing at B1.
\end{grammarlens}

\subsection*{Your turn}
\begin{itemize}
\raggedright
  \item \emph{Say it:} \emph{il portico}
  \item \emph{Say it:} \emph{il portico}, once more
  \item \emph{Say it:} a sentence of your own with \emph{il portico}
  \item \emph{Recall:} say the Italian for the main church, then the Italian for the cathedral, then say \emph{il portico} again
\end{itemize}

\begin{tcolorbox}[breakable,colback=teal!4,colframe=teal!35!black,title={Before you move on}]
What does il portico mean? (\textquotedblleft{}The arcade, the portico\textquotedblright{}.) Say it once more.
\end{tcolorbox}
\section[la loggia]{la loggia — the loggia, the open gallery}

\label{lesson:IT-C445-loggia}



\begin{quote}
Before the new one: say the Italian for the cathedral, then the Italian for the arcade.
\end{quote}

\subsection*{You'll want to know: la loggia}
\textbf{la loggia} — \textquotedblleft{}the loggia, the open gallery\textquotedblright{}.

\begin{quote}
Dalla loggia del palazzo si vede la piazza.
\end{quote}

\begin{quote}
\emph{The square is in view from the loggia of the palace.}
\end{quote}

A covered space open on one side.

\begin{cousinweb}
\emph{loggia} was borrowed from Old French \emph{loge} 'arbour, covered walk' (EtymDB), from Frankish \emph{laubja} 'arbour'.
\end{cousinweb}

\begin{grammarlens}[title={What you've built}]
One more word for telling, explaining and describing at B1.
\end{grammarlens}

\subsection*{Your turn}
\begin{itemize}
\raggedright
  \item \emph{Say it:} \emph{la loggia}
  \item \emph{Say it:} \emph{la loggia}, once more
  \item \emph{Say it:} a sentence of your own with \emph{la loggia}
  \item \emph{Recall:} say the Italian for the cathedral, then the Italian for the arcade, then say \emph{la loggia} again
\end{itemize}

\begin{tcolorbox}[breakable,colback=teal!4,colframe=teal!35!black,title={Before you move on}]
What does la loggia mean? (\textquotedblleft{}The loggia, the open gallery\textquotedblright{}.) Say it once more.
\end{tcolorbox}
